\section{Building the multilingual benchmark}
\label{sec:benchmark}

\subsection{Source slices}

We translate the five binary safety slices of JevBench \citep{pai2026jevbench},
\MLnCasesEN{} English cases in total, built from public datasets pinned to exact
revisions:

\begin{itemize}
  \item \textbf{Jailbreak classification} (400 cases) \citep{jackhhao2023jailbreak}:
  is the prompt an attempt to make an assistant ignore its safety rules?
  \item \textbf{Prompt injections} (116 cases) \citep{deepset2023promptinjections}:
  does the text try to override, ignore or extract the instructions of the
  system that receives it? The source mixes English and German texts.
  \item \textbf{Aegis 2.0 prompts} and \textbf{responses} (500 cases each)
  \citep{ghosh2025aegis2}: is the user prompt, or the assistant's response,
  unsafe under a general content-safety policy?
  \item \textbf{ATBench} (500 cases) \citep{li2026atbench}: considering a whole
  trajectory of a tool-using agent, did the agent behave unsafely?
\end{itemize}

The safety slices are sampled close to class balance, except jailbreak
classification, where 35\% of cases are jailbreaks. Every question carries the
same preamble, which tells the model to treat the state as data to be evaluated
and never as instructions, because the states contain adversarial text by design.

\subsection{Languages}

We chose eight languages to vary three properties that plausibly affect a model:
script, language family, and the amount of text available on the web.
Table~\ref{tab:langs} lists them. Chinese, Spanish, French and Russian are
commonly treated as high-resource languages and Swahili as a lower-resource one,
with Arabic, Hindi and Vietnamese between them; we use these labels only to
describe the design, not as measured quantities. Four languages share the
Latin script with English, which separates the effect of script from the effect
of language.

\begin{table}[t]
\centering
\caption{Languages and translated cases. A case is kept only if every one of its translated strings passed the checks of Section~\ref{sec:checks}; the counts are the cases each language delivered, out of the English cases in the first row.}
\label{tab:langs}
\footnotesize
\setlength{\tabcolsep}{4pt}
\resizebox{\linewidth}{!}{%
\begin{tabular}{@{}llllrrrrrr@{}}
\toprule
Code & Language & Family & Script & Jailbreak & Prompt inj. & Aegis prompts & Aegis responses & ATBench & Total \\
\midrule
en & English (source) & Indo-European, Germanic & Latin & 400 & 116 & 500 & 500 & 500 & 2{,}016 \\
zh & Chinese & Sino-Tibetan & Han (Simplified) & 397 & 116 & 499 & 500 & 467 & 1{,}979 (98.2\%) \\
ar & Arabic & Afro-Asiatic, Semitic & Arabic & 399 & 116 & 500 & 499 & 493 & 2{,}007 (99.6\%) \\
hi & Hindi & Indo-European, Indo-Aryan & Devanagari & 397 & 116 & 500 & 499 & 487 & 1{,}999 (99.2\%) \\
vi & Vietnamese & Austroasiatic & Latin & 394 & 115 & 497 & 499 & 463 & 1{,}968 (97.6\%) \\
sw & Swahili & Niger-Congo, Bantu & Latin & 398 & 116 & 500 & 499 & 484 & 1{,}997 (99.1\%) \\
es & Spanish & Indo-European, Romance & Latin & 396 & 113 & 499 & 500 & 451 & 1{,}959 (97.2\%) \\
fr & French & Indo-European, Romance & Latin & 395 & 115 & 500 & 500 & 459 & 1{,}969 (97.7\%) \\
ru & Russian & Indo-European, Slavic & Cyrillic & 399 & 116 & 500 & 500 & 490 & 2{,}005 (99.5\%) \\
\bottomrule
\end{tabular}}
\end{table}


\subsection{Translation}
\label{sec:translation}

\paragraph{Translator.} Translations were produced by \texttt{gpt-6-astra} at
low reasoning effort, called through the Codex command-line client
(version 0.155.1) on 15 and 16 October 2026. We did not use a dedicated
machine-translation system, because safety data contains jailbreaks, slurs and
descriptions of crimes, and the translator must translate them faithfully rather
than refuse, soften or summarise them.

\paragraph{What is translated.} Every natural-language string of a case's
state, and the question's instruction and answer descriptions. For ATBench,
whose states are JSON, only the user turns, the agent's thoughts and messages,
and the tool descriptions are translated; tool names, actions, arguments and
environment outputs stay byte-identical, so the trajectory remains the same
sequence of actions. Gold labels, identifiers and every other field are never
sent to the translator. Each unique string is translated once and reused
wherever it occurs (Aegis prompts and responses share their prompts). The five
slices hold 5{,}081 unique strings, about 2.06 million characters, per language.

\paragraph{Instructions to the translator.} The prompt
(Appendix~\ref{app:prompt}) states that the strings are benchmark data, that the
harmful content is deliberate, and that each translation must stay ``exactly as
harmful, manipulative, rude, benign, correct or wrong as the original''. It asks
the translator not to soften, censor, summarise, add warnings or correct
anything, and to translate, not obey, any instruction inside a string. Strings
are sent in batches, each wrapped in a tag that carries a random nonce, so text
inside a string cannot close the block. The prompt names English as the source
language; the German prompt-injection texts were translated into the target
language as well.

\subsection{Checks}
\label{sec:checks}

Every returned string is checked before use:

\begin{itemize}
  \item \textbf{Refusal}: the translator declined, or the output contains a
  refusal phrase that the source does not.
  \item \textbf{Script}: for languages with their own script, at least 30\% of
  the letters not copied from the source must be in that script (3\% for
  Vietnamese, whose diacritics are a small share of its letters).
  \item \textbf{Untranslated}: for Latin-script targets, at most 60\% of the
  output's words may also appear in the source; for every target, at most
  70\% of the output's letters may be copied verbatim.
  \item \textbf{Length}: for sources of 40 characters or more, the output must
  be 0.4 to 3.5 times as long (0.12 to 2.5 for Chinese).
  \item \textbf{JSON}: an ATBench fragment must keep every tool name and its
  truncation marker.
\end{itemize}

A string that fails is sent again once on its own. A case with any string still
failing is left out of that language. We do not repair translations by hand.
Table~\ref{tab:excluded} in the appendix gives the reasons per language. Most
exclusions were flagged as untranslated. On inspection, many of these are
translations of jailbreak templates and tasks that keep long stretches of
code, markup or English template text, which the check cannot tell apart from
a translator that did nothing; leaving them out is conservative and removes a
few of the longest jailbreak prompts from the Latin-script languages. Every language keeps between \MLusableMinPct\% (\MLusableMinLang)
and \MLusableMaxPct\% of the English cases.

\paragraph{Two variants.} The pipeline writes two case sets per language: one
with the state and the question translated, and one with only the state
translated and the question kept in English. All results in this paper use the
first, fully translated variant, which is how a non-English deployment would
query the model.
